% Part: history
% Chapter: biographies 
% Section: georg-cantor
%READERNOTE{SOL6-A272: कँटर यांच्या उपचाराच्या आरंभाचे स्रोतामधील 1894 हे वर्ष 1884 केले आहे. ऑलिव्हर डाइझर यांच्या Einführung in die Mengenlehre मधील चरित्रकालक्रम, पृ. 218, मे 1884 मध्ये हाल येथील आरोग्यधामातील वास्तव्य नोंदवतो; लिओपोल्डिनाचे सदस्यचरित्रही 1884 पासून आजार आणि पुन्हा पुन्हा उपचार नोंदवते: https://www.leopoldina.org/en/members/member-list/detail/georg-cantor. येथे संपूर्ण वैद्यकीय नोंदींचे स्वतंत्र पुनरावलोकन केलेले नाही.}

\documentclass[../../../include/open-logic-section]{subfiles}

\begin{document}

\olfileid{his}{bio}{can}

\olsection{गेऑर्ग कँटर}

\olphoto{cantor-georg}{गेऑर्ग कँटर}

गेऑर्ग कँटर (\textsc{गे}-ऑर्ग \textsc{कान}-टोर) यांच्या एका जुन्या चरित्रात असा दावा केला होता की रशियातील सेंट पीटर्सबर्गकडे जाणाऱ्या जहाजावर त्यांचा जन्म झाला आणि ते तिथे सापडले; त्यांचे पालक कोण होते हे माहीत नव्हते. मात्र हा दावा खरा नाही. कँटर यांचा जन्म १८४५ मध्ये सेंट पीटर्सबर्ग येथे झाला होता.

कँटर यांनी १८६७ मध्ये बर्लिन विद्यापीठातून गणितातील डॉक्टरेट मिळवली. संच उपपत्तीतील त्यांच्या कार्यासाठी ते ओळखले जातात आणि संशोधनाची स्वतंत्र शाखा म्हणून संच उपपत्तीची स्थापना करण्याचे श्रेय त्यांना दिले जाते. निरनिराळ्या आकारमानांचे अनंत संच असतात हे सर्वप्रथम त्यांनी सिद्ध केले. त्यांच्या उपपत्तींबद्दल, विशेषतः अनंतांच्या उपपत्तीबद्दल, तत्कालीन गणितज्ञांत बरीच चर्चा झाली आणि त्यांचे कार्य वादग्रस्त ठरले.

कँटर यांच्या धार्मिक श्रद्धा आणि गणितीय कार्य यांची घट्ट सांगड होती. अनंतपार संख्यांची उपपत्ती आपल्याला थेट ईश्वराकडून कळली, असा दावाही त्यांनी केला होता. आयुष्याच्या पुढील काळात त्यांना मानसिक आजाराचा त्रास झाला. १८८४ पासून त्यांना रुग्णालयात दाखल करावे लागले; नंतरच्या काळात असे अधिक वारंवार घडले. या स्रोताच्या मांडणीनुसार, त्यांच्या कार्यावरील तीव्र टीका आणि गणितज्ञ लिओपोल्ड क्रोनेकर यांच्याशी झालेला मतभेद यांमुळे त्यांना नैराश्य आले आणि गणितातील रस कमी झाला. नैराश्याच्या काळात कँटर तत्त्वज्ञान आणि साहित्याकडे वळत. शेक्सपियरची नाटके फ्रान्सिस बेकन यांनी लिहिली होती, असा सिद्धांतही त्यांनी प्रकाशित केला.

६ जानेवारी १९१८ रोजी हाल येथील आरोग्यधामात कँटर यांचे निधन झाले.

\begin{reading}
कँटर यांच्या सविस्तर चरित्रांसाठी \citet{Dauben1990} आणि \citet{Grattan-Guinness1971} पाहा. त्यांच्या अपरंपरागत मतांचे वर्णन बीबीसी रेडिओ ४ वरील \emph{A Brief History of Mathematics} \citep{Sautoy2014} या कार्यक्रमातही आहे. कँटर यांच्या उपपत्ती रॅपच्या रूपात ऐकायच्या असतील, तर \citet{Rose2012} पाहा.
\end{reading}

\end{document}
